%% ortak_baslik.tex -- iki iskelette de aynı: paketler, \sonuc makrosu, örnek uyarısı.
\usepackage[T1]{fontenc}
\usepackage{amsmath}
\usepackage{graphicx}
\graphicspath{{sekiller/}{./}}   %% şekil sekiller/ klasöründe ya da kökte (düz gönderim zip'i) bulunur
\usepackage{booktabs}
\usepackage{tabularx}
\usepackage{algorithm}
\usepackage{algpseudocode}
\usepackage{tikz}
\usetikzlibrary{positioning}
\usepackage{url}

%% --- Sonuç makroları: \sonuc{anahtar}. Değer yoksa PDF'te kalın "??anahtar" çıkar (sessiz geçmez). ---
\makeatletter
\newcommand{\sonucyaz}[2]{\expandafter\gdef\csname sonuc@#1\endcsname{#2}}
\newcommand{\sonuc}[1]{\@ifundefined{sonuc@#1}{\textbf{??\detokenize{#1}}}{\csname sonuc@#1\endcsname}}
\makeatother
% OTOMATİK ÜRETİLDİ (sonuc_uret.py) -- elle değiştirmeyiniz; sonuclar.csv değişince yeniden üretiniz.
\sonucyaz{kat.n}{10}
\sonucyaz{base.ort}{0.704}
\sonucyaz{base.ss}{0.016}
\sonucyaz{yeni.ort}{0.770}
\sonucyaz{yeni.ss}{0.023}
\sonucyaz{fark.ort}{0.066}
\sonucyaz{fark.alt}{0.054}
\sonucyaz{fark.ust}{0.078}
\sonucyaz{fark.p}{= 0.002}
\sonucyaz{fark.ptest}{exact}
\sonucyaz{fark.npos}{10}
\sonucyaz{yayin.ort}{0.773}
\sonucyaz{yayin.ss}{0.019}
\sonucyaz{yfark.ort}{\ensuremath{-}0.003}
\sonucyaz{yfark.alt}{\ensuremath{-}0.008}
\sonucyaz{yfark.ust}{0.002}
\sonucyaz{yfark.p}{= 0.256}
\sonucyaz{yfark.ptest}{exact}
\sonucyaz{yfark.npos}{4}
\sonucyaz{veri.kayit}{400}
\sonucyaz{veri.kisi}{100}
\sonucyaz{veri.fs}{2000}
   %% sonuc_uret.py üretir (elle doldurulmaz); başka yerde sonuç sayısı yazılmaz

%% --- Örnek uyarısı: kendi metninizi yazınca govde.tex'in başındaki \ornekuyarisi satırını siliniz. ---
\newcommand{\ornekuyarisi}{\par\noindent\fbox{\parbox{\dimexpr\linewidth-2\fboxsep-2\fboxrule}{\footnotesize
\textbf{ÖRNEK İSKELET.} Bu belgedeki cümleler ve sayılar yalnız biçimi göstermek içindir; gerçek
bir çalışmanın sonucu değildir. Kendi verinizle değiştiriniz, sonra bu kutuyu siliniz.}}\par\medskip}
